\chapter{BIM and Alignment Modelling}

\section{Motivation}

\begin{remark}
Building Information Modelling (BIM) has become a central paradigm in
modern infrastructure projects.
\end{remark}

\begin{remark}
Its goal is the integration of geometry, semantics, and lifecycle
information into a unified digital representation.
\end{remark}

\begin{remark}
However, despite its broad scope, the representation of railway
alignments within BIM remains conceptually and technically limited.
\end{remark}

---

\section{BIM as a Data Integration Paradigm}

\begin{remark}
BIM is fundamentally concerned with:
\begin{itemize}
\item object-based modelling,
\item semantic enrichment of geometry,
\item interoperability across disciplines,
\item lifecycle management of infrastructure assets.
\end{itemize}
\end{remark}

\begin{remark}
In this context, geometry is typically associated with objects such as:
\begin{itemize}
\item buildings,
\item bridges,
\item track elements,
\item terrain models.
\end{itemize}
\end{remark}

\begin{remark}
Alignment geometry is often embedded implicitly within these objects,
rather than treated as a primary modelling entity.
\end{remark}

---

\section{Alignment in BIM Standards}

\subsection{IFC Alignment}

\begin{remark}
Recent extensions of IFC introduce explicit alignment entities.
\end{remark}

\begin{remark}
These aim to represent:
\begin{itemize}
\item horizontal alignment,
\item vertical alignment,
\item cant (superelevation),
\item stationing.
\end{itemize}
\end{remark}

\begin{remark}
While this is a significant step forward, several issues remain:
\begin{itemize}
\item separation of components,
\item lack of unified state representation,
\item limited support for dynamic interpretation,
\item complexity of schema usage.
\end{itemize}
\end{remark}

---

\subsection{Layered Representation}

\begin{remark}
In most BIM-related formats, alignment is decomposed into layers:
\begin{itemize}
\item horizontal geometry,
\item vertical profile,
\item cant definition.
\end{itemize}
\end{remark}

\begin{remark}
This decomposition is practical for data exchange, but does not directly
support computation or optimization.
\end{remark}

---

\section{Limitations of Current BIM Alignment Modelling}

\subsection{Static representation}

\begin{remark}
BIM primarily represents geometry as static data.
\end{remark}

\begin{remark}
Dynamic properties such as acceleration, jerk, or vehicle response are
not part of the core model.
\end{remark}

---

\subsection{Fragmentation}

\begin{remark}
The separation of horizontal, vertical, and cant components leads to a
fragmented representation.
\end{remark}

\begin{remark}
This fragmentation complicates:
\begin{itemize}
\item consistency checks,
\item dynamic evaluation,
\item optimization-based design.
\end{itemize}
\end{remark}

---

\subsection{Lack of computational core}

\begin{remark}
BIM standards focus on data exchange and documentation.
\end{remark}

\begin{remark}
They do not provide a unified computational model for alignment
analysis and optimization.
\end{remark}

---

\section{AIM as a Computational Layer for BIM}

\subsection{Conceptual Role}

\begin{remark}
The AIM framework can be understood as a computational core that
complements BIM data models.
\end{remark}

\begin{remark}
Instead of replacing BIM, AIM operates between:
\begin{itemize}
\item external data formats,
\item and internal engineering reasoning.
\end{itemize}
\end{remark}

---

\subsection{Canonical Representation}

\begin{remark}
AIM introduces a unified representation:
\[
(\gamma(s), \theta(s), \kappa(s), \phi(s)).
\]
\end{remark}

\begin{remark}
This representation:
\begin{itemize}
\item integrates horizontal geometry,
\item integrates vertical and cant behaviour,
\item supports dynamic evaluation,
\item is directly usable in optimization.
\end{itemize}
\end{remark}

---

\subsection{Transformation Pipeline}

\begin{verbatim}
BIM / IFC / LandXML
→ parsing
→ fat representation
→ normalization
→ AIM state model
→ analysis / optimization
→ export
\end{verbatim}

\begin{remark}
This pipeline separates:
\begin{itemize}
\item data exchange concerns,
\item computational modelling,
\item engineering reasoning.
\end{itemize}
\end{remark}

---

\section{Alignment-Centred BIM}

\subsection{Shift of perspective}

\begin{remark}
The AIM framework suggests a shift toward alignment-centred modelling.
\end{remark}

\begin{remark}
Instead of treating alignment as one component among many, it becomes a
primary organizing structure.
\end{remark}

---

\subsection{Implications}

\begin{itemize}
\item infrastructure elements are referenced to alignment,
\item geometry is derived from alignment state,
\item dynamic behaviour becomes part of the model,
\item optimization becomes a native operation.
\end{itemize}

---

\section{Integration with BIM Workflows}

\subsection{Import}

\begin{remark}
Existing BIM data is imported via:
\begin{itemize}
\item IFC parsing,
\item LandXML parsing,
\item legacy format interpretation.
\end{itemize}
\end{remark}

---

\subsection{Processing}

\begin{remark}
Within AIM, the imported data is:
\begin{itemize}
\item normalized,
\item reconstructed,
\item enriched with dynamic interpretation,
\item potentially optimized.
\end{itemize}
\end{remark}

---

\subsection{Export}

\begin{remark}
Results may be exported back into:
\begin{itemize}
\item IFC alignment structures,
\item LandXML files,
\item other engineering formats.
\end{itemize}
\end{remark}

\begin{remark}
This ensures compatibility with existing BIM ecosystems.
\end{remark}

---

\section{Relation to BIMinfra}

\begin{remark}
Infrastructure BIM (often referred to as BIMinfra) extends BIM concepts
to linear infrastructure such as railways and roads.
\end{remark}

\begin{remark}
AIM can be interpreted as a specialization within BIMinfra focusing on:
\begin{itemize}
\item alignment-centric modelling,
\item dynamic interpretation,
\item optimization-driven design.
\end{itemize}
\end{remark}

---

\section{Discussion}

\begin{remark}
BIM provides a powerful framework for data integration, but lacks a
strong computational alignment model.
\end{remark}

\begin{remark}
AIM addresses this gap by:
\begin{itemize}
\item introducing a unified state representation,
\item enabling dynamic evaluation,
\item supporting optimization,
\item and integrating heterogeneous data sources.
\end{itemize}
\end{remark}

\begin{remark}
Together, BIM and AIM can be seen as complementary:
\begin{itemize}
\item BIM for data integration and lifecycle,
\item AIM for computation and alignment reasoning.
\end{itemize}
\end{remark}

---

\section{Summary}

\begin{summary}
Current BIM standards provide important structures for data exchange and
semantic modelling, but remain limited in their treatment of railway
alignment.

The AIM framework complements BIM by introducing a unified,
state-based, and optimization-ready representation of alignment.

This enables a transition from static data handling toward dynamic,
computational, and alignment-centred infrastructure modelling.
\end{summary}
